\bookchapter{Chapter 0: Hi, Before We Begin}

Before we officially start exploring, let's cover some basic rules to make sure you understand everything that follows. Don't rush—sharpening the axe won't delay the woodcutting.

\booksection{What This Book Can Do for You (and What It Can't)}

It can help you systematically understand how to use your phone, grasp the reasoning behind features, and—whether you're an Android fan or an Apple devotee—most of the thinking applies universally. It will guide you through hands-on practice, taking you from “able to use” to “truly understanding how to use.”

It is not a manual for any specific phone model, nor is it tech news chasing the latest gadgets, and it's certainly not a guide to becoming addicted to your phone. Honestly, if you're expecting “how to enjoy scrolling videos more,” you may be disappointed—we care far more about making your phone truly serve you, not the other way around.

\booksection{Our “Secret Code” Guide}

To be clearer, the book uses some small markers—these markers are like tacit codes between us:

\begin{itemize}
  \item \textbf{[Button Name]}: Refers to a specific button or App name on the screen, such as [Settings], [OK].
  \item \textbf{>}: Indicates the order of operations. For example, "Tap [Settings] > [Sound]" means first tap Settings, then tap Sound.
  \item \textbf{Bold}: Very important! It may be a key concept or a safety warning—please pay close attention.
  \item \emph{Italics}: Supplementary thoughts or analogies to help you understand better.
\end{itemize}

\booksection{Android and iOS—We Cover Both}

Guess what—some things have different names but do the same thing. For example: Application (Android) and App (iOS), Settings (called Settings on both), File Manager (Android) and Files (iOS). Knowing they're the same thing is enough; don't get tripped up by the terminology.

Additionally, within the Android camp there's Huawei's HarmonyOS, which has its own characteristics—we'll cover it specifically later. In short, no matter what phone you use, this book can help you.

\booksection{Phone Storage: Why Did You Buy 128GB but Only Have Less Than 110GB Available?}

Everyone knows the storage units: 1TB = 1024GB, 1GB = 1024MB, 1MB = 1024KB... But here's the catch: hard drive manufacturers usually calculate in base 1000, while the system calculates in base 1024, so the “128G” drive you bought will “shrink” in actual capacity.

But there are three main reasons your phone's storage shrinks:

\begin{enumerate}
  \item \textbf{The system itself takes up space}: The operating system is like a house's load-bearing walls—it must exist. Modern systems routinely occupy 10-20GB, which is completely normal.
  \item \textbf{Manufacturer pre-installed apps}: The apps that come with a new phone also take up space; some can be uninstalled, some cannot.
  \item \textbf{Partitioning and formatting losses}: For technical reasons—just like the lining of a backpack also takes up volume.
\end{enumerate}

\begin{tcbinfo}{Tips for Buying a Phone}
Nowadays 128GB phones are basically not enough; 256GB is a safer starting point. If you like taking photos or playing games, consider going straight to 512GB or larger. Remember—when it comes to storage, it's better to have surplus than to come up short.

\end{tcbinfo}

\begin{tcbpractice}{Hands-on: Check Your Phone's "ID Card"}
Open [Settings], find About Phone (or a similar option), and note down: phone model, system version, total storage, and available storage. This is the first step to understanding it. Also check battery health so you have a clear picture.

\end{tcbpractice}

\booksection{Your Phone Is a “Digital Asset Vault”}

Before diving into specific knowledge, let's establish a core understanding:\textbf{your phone has long been more than a communication tool—it is your most important digital asset vault}. Contacts, chat history, family photos, work documents, banking apps, payment codes, health data, electronic IDs... Together, this data is often worth more than the phone itself.

\begin{tcbinfo}{Why Does This Matter?}
If your phone is lost or broken, you lose not only a few thousand yuan of hardware,\textbf{but more likely years of irreplaceable data}—photos of your child growing up, voice messages from departed loved ones, important contracts that were never backed up. Scientific organization and protection habits can safeguard you across three dimensions:

\begin{itemize}
  \item \textbf{Improve operating efficiency}: Standardized file classification makes system lookups faster and keeps the phone from lagging over the long term.
  \item \textbf{Extend service life}: Reasonable storage management reduces the read/write burden on the chip, making both battery and storage last longer.
  \item \textbf{Ensure data security}: Layered management + regular backups + tightened permissions keep "data disasters" away from you.
\end{itemize}

\end{tcbinfo}

\booksubsection{Our Core Philosophy: “Prevention First, Combined with Treatment”}

Just as you can't wait until you're sick to see a doctor for wellness,\textbf{phone management must start with daily habits}. What this book repeatedly emphasizes is not some advanced trick, but these simple principles:

\begin{itemize}
  \item \textbf{Prevention first}: Install apps only from official stores, strip metadata from photos before sharing, never reuse passwords—these "small actions" cost very little but block the vast majority of risks.
  \item \textbf{Regular check-ups}: Organize the download folder once a month, audit permissions once a quarter, do a full backup every six months—a fixed rhythm is a hundred times more effective than ad-hoc firefighting.
  \item \textbf{Always have a fallback}: Enable cloud sync, keep account passwords recorded, write a "digital estate checklist"—so that when accidents happen, you and your family can respond calmly.
\end{itemize}
